% ---
% filename : oibt_securite_20260517_000_cinq_regles_securite.tex
% id : oibt_securite_20260517_000
% title : "Les cinq règles de sécurité pour les travaux hors tension"
% domain : "OIBT"
% subdomain : "Sécurité"
% tags : ["oibt", "sécurité", "mise_hors_tension", "travaux_électriques", "dt"]
% difficulty : 1
% time_solve : 8
% has_octave : false
% author : "om"
% ---
\documentclass[a4paper,11pt,answers,addpoints]{exam}
% ---------------------------------------------------------
% LANGUE, ENCODAGE, FONTS
% ---------------------------------------------------------
\usepackage[utf8]{inputenc}
% ---------------------------------------------------------
% GÉOMÉTRIE ET MISE EN PAGE
% ---------------------------------------------------------
\usepackage[left=1.7cm,right=1.5cm,top=2cm,bottom=1cm]{geometry}
\usepackage{microtype}
\usepackage{float}
\usepackage{needspace}

% ---------------------------------------------------------
% MATHÉMATIQUES
% ---------------------------------------------------------
\usepackage{amsmath, amssymb, bm, esvect}

\newcommand{\V}{\overrightarrow}
\def\vect#1{\overrightarrow{\strut#1}}
\providecommand{\Degrees}[0]{\ensuremath{^\circ}}

% ---------------------------------------------------------
% UNITÉS ET GRANDEURS (SIunitx)
% ---------------------------------------------------------
\usepackage{siunitx}
\sisetup{output-decimal-marker={,}}
\DeclareSIUnit \var {var}
\DeclareSIUnit \VA {VA}
\DeclareSIUnit \kVA {kVA}
\DeclareSIUnit[quantity-product={}] \degree{\text{\textdegree}}

% Permet la césure dans les nombres avec unités (évite les Underfull \hbox)
%\AllowBreakAtQQ

% ---------------------------------------------------------
% GRAPHIQUES : TikZ + CircuitikZ (sans tkz-fct qui cause des erreurs spath3)
% ---------------------------------------------------------
\usepackage{tikz}
\usetikzlibrary{
	arrows.meta,
	intersections,
	decorations.pathreplacing,
	positioning
}

\usepackage[european,straightvoltages,EFvoltages]{circuitikz}
\usepackage{pgfplots}
\pgfplotsset{compat=1.12}

% Symbole étoile-triangle personnalisé
\newcommand{\wye}{%
	\mathbin{\tikz[x=1ex,y=1ex]{%
			\draw[line width=.1ex] (0,0)--(45:1)--++(-45:1) (45:1)--++(0,1);
	}}%
}

% ---------------------------------------------------------
% TABLEAUX
% ---------------------------------------------------------
\usepackage{tabularx}
\usepackage{tabu}

% ---------------------------------------------------------
% HYPERLIENS
% ---------------------------------------------------------
\usepackage[colorlinks=true,
menucolor=red,
breaklinks=true,
urlcolor=blue,
linkcolor=blue]{hyperref}

% ---------------------------------------------------------
% COMMANDES PERSONNELLES
% ---------------------------------------------------------
\newcommand\nomauteur{bdminasmorgul@protonmail.com}
\newcommand\dateTe{18.05.2026}
\newcommand\brancheTe{DT}

% ---------------------------------------------------------
% STYLE DE L'EXERCICE
% ---------------------------------------------------------
\pagestyle{headandfoot}
\firstpageheadrule
\runningheadrule

% En-tête avec largeur fixe pour éviter les dépassements
\firstpageheader{\brancheTe\ (\nomauteur)}{Élève : }{\makebox[3.5cm][r]{\dateTe\ \thepage/\numpages}}
\runningheader{\brancheTe\ (\nomauteur)}{Élève : }{\makebox[3.5cm][r]{\dateTe\ \thepage/\numpages}}

% Pied de page
\newcommand{\totalpointtts}{%
	\begin{tikzpicture}[remember picture, overlay]
		\node[anchor=south west, inner sep=0pt, outer sep=0pt]
		at ([xshift=0.5cm,yshift=0.5cm]current page.south west)
		{\rotatebox{90}{Total des points : \numpoints}};
	\end{tikzpicture}%
}

\firstpagefooter{\hspace*{\fill}\totalpointtts}{}{}
\runningfooter{\hspace*{\fill}\totalpointtts}{}{}

% Réglages exam
\printanswers
\addpoints
\pointsinmargin
\bracketedpoints

% ---------------------------------------------------------
% LISTINGS (code)
% ---------------------------------------------------------
\usepackage{xcolor}
\usepackage{listings}

\lstdefinestyle{octavestyle}{
	language=Matlab,
	basicstyle=\ttfamily\small,
	keywordstyle=\color{blue},
	commentstyle=\color{green!50!black},
	stringstyle=\color{red},
	stepnumber=1,
	frame=single,
	breaklines=true,
	breakatwhitespace=true,
	tabsize=4
}

\usepackage{enumitem}
\setlist[enumerate,1]{label=\alph*), ref=\alph*)}
\newenvironment{explication}{%
	\par\noindent\textbf{Explication. }\itshape
}{\par\normalfont}

\begin{document}
	\unframedsolutions
	\pointsinmargin
	\bracketedpoints
	
	\section*{Préparation DT PQ CFC ELMO,IELE,PELE}
	\begin{questions}
\question
En tant qu'installateur-électricien, vous devez effectuer le remplacement d'un départ moteur dans une armoire industrielle. Pour garantir votre sécurité et celle de vos collègues, vous devez impérativement appliquer les mesures de protection contre les chocs électriques lors de travaux hors tension conformément à l'OIBT.

\begin{parts}
	\part[3] Une fois la partie de l'installation identifiée, quelle est la toute première opération physique à effectuer sur l'organe de coupure~?
	\part[3] Après avoir ouvert le circuit, quelle mesure devez-vous prendre pour empêcher qu'une tierce personne ne rétablisse le courant par inadvertance~?
	\part[3] Avant tout contact direct avec les parties actives, quel test est absolument indispensable de réaliser avec un appareil de mesure adapté~?
	\part[3] Dans quelle condition spécifique est-il obligatoire de mettre l'installation en court-circuit et à la terre~?
\end{parts}

\begin{solution}
	\begin{itemize}
		\item \textbf{Déclencher} : Selon l'OIBT Art. 22 let. a, il faut séparer l'installation des sources d'énergie (couper les disjoncteurs ou sectionneurs).
		\item \textbf{Sécurisation} : Selon l'OIBT Art. 22 let. b, assurer contre le réenclenchement (verrouillage par cadenas ou signalisation claire).
		\item \textbf{Vérification} : Selon l'OIBT Art. 22 let. c, vérifier l'absence de tension sur tous les pôles à l'aide d'un testeur de tension (DDT/VAT).
		\item \textbf{Mise à terre} : Selon l'OIBT Art. 22 let. d, mettre en court-circuit et à la terre s'il existe un danger de tension induite (proximité d'autres lignes) ou de retour de tension (génératrices, photovoltaïque).
		\item \textbf{Protection} : Selon l'OIBT Art. 22 let. e, protéger les parties voisines restées sous tension par des écrans ou des protecteurs isolants.
	\end{itemize}
	
	\begin{align*}
		\text{Étape 1} &= \text{Déclencher} \\[6pt]
		\text{Étape 2} &= \text{Assurer contre le réenclenchement} \\[6pt]
		\text{Étape 3} &= \text{Vérifier l'absence de tension} \\[6pt]
		\text{Étape 4} &= \text{Mettre en court-circuit et à la terre (si danger)} \\[6pt]
		\text{Étape 5} &= \text{Protéger des parties voisines}
	\end{align*}
\end{solution}
\end{questions}
\end{document}
